\paragraph{Recurrence, later exclusions, and valuations.}
Only now impose the extra finite coprimality sieve $B$.
The zero-or-retain Fourier sieve is a finite additive translation
operator. It preserves cuspidality, keeps every coefficient whose
index is prime to $B$, and commutes with good $T(p^2)$ for
$p\notin S_1=S_0\cup B$. In particular it keeps $a(D_0)=1$.
The squarefree comparison (E.11) was proved before this enlargement,
so it requires no extra squareclass hypotheses at the primes of $B$.
After the sieve it is asserted at the retained indices.
The standard good $T(p^2)$ formula, with $\psi^2=\eta^2$, gives
\[
 a(p^2n)=\{\eta(p)a_p(E)-\psi(p)(-n/p)\}a(n)
                       -p\eta(p)^2a(n/p^2),\qquad p\notin S_1.
                                                        \tag{E.12}
\]
For squarefree $D$, write $a(Dp^{2r})=\eta(p)^r b_r$.
The first multiplier is
$a_p(E)-(\psi/\eta)(p)(-D/p)$; subsequent multipliers satisfy
$b_{r+1}=a_p(E)b_r-pb_{r-1}$. All are integers, including when
$p\mid D$. Iteration over the good primes gives the same assertion
for every retained square shell after removal of the $\eta$ phase.
More explicitly, for $p\notin S_1$ set
$A_{p,0}(D)=1$, $A_{p,1}(D)=a_p(E)-(\psi/\eta)(p)(-D/p)$,
and
$A_{p,r+1}(D)=a_p(E)A_{p,r}(D)-pA_{p,r-1}(D)$ for $r\ge1$.
Then, for every retained index $Dm^2$ with $D$ squarefree,
\[
 a(Dm^2)=z_D\eta(m)r_D
              \prod_{p^r\Vert m}A_{p,r}(D).
\]
The formula includes $p\mid D$, when the initial Legendre symbol
is zero. When applying the recurrence at one prime, the other
prime-square factors do not alter that Legendre symbol. This proves
the formula by induction on the prime factors of $m$, with
$a(n/p^2)=0$ when $p^2\nmid n$. All factors $A_{p,r}(D)$ are
integers. Thus every coefficient used by the extraction and shell
operators of Proposition~\ref{ex:app-coefficient-rules} is a rational
number times a two-power root of unity, including coefficients with
non-squarefree index. In the chosen residue field such a root reduces
to one. A rational coefficient of valuation zero reduces to the
element $1$ of its prime field $\F_2$. Consequently the fresh
coefficient-one reading of that proposition detects precisely equality
in the fixed minimum-depth bound. It uses the same $g$, $D_0$ and
normalization throughout.


Combining (E.9) and (E.11) makes the values
$v_2a(D)-\omega(D)$ at nonzero retained squarefree coefficients a
bounded-below subset of $\Z$. It is nonempty because $D_0$ is retained,
and therefore has an attained finite minimum. The form and the initial
finite sieve are fixed before this minimum is chosen. Restricting
later fresh primes by further exclusions does not authorize choosing
another minimum.
\end{proof}

% End 07-even-source.tex


